% Owner: Hassan Ashraf 24L-3073
\section{Conclusion}
\label{sec:conclusion}

Deliverable~1 established the foundations of ProcurePulse. We defined the
problem: small and medium-sized businesses that buy metals face large and
unpredictable price movements, and cannot easily judge whether a supplier's
quote is fair. We built a reproducible pipeline that collects 16 years of daily
data for four metals and four external signals, cleans it and turns it into
model-ready features (Sections~\ref{sec:collection} to~\ref{sec:features_eda}).

Deliverable~2 trained and evaluated the first models. We framed the task as
predicting whether the gold price will be higher in five trading days, and
compared a majority baseline, logistic regression, a random forest and
histogram gradient boosting (Section~\ref{sec:classical}) with an LSTM
(Section~\ref{sec:deep}). All models were tuned with walk-forward
cross-validation on the training years, selected on the validation years and
tested once on 2024--2026 (Section~\ref{sec:evaluation}).

The results are modest and informative. The best test AUC is 0.60, so the
direction of gold is only slightly predictable from these features. Logistic
regression produced the most reliable ``up'' signal, with 70\% precision on
its ``up'' calls, and the LSTM, which uses only stationary features, gave the
most balanced predictions (F1 0.55). The tree models ranked days better than
chance but predicted ``down'' on every test day, because they rely on price
levels that moved far outside the training range. Accuracy alone would have
hidden all of this, which is why we report precision, recall, F1 and AUC-ROC
and analyse errors by market regime (Section~\ref{sec:error-bias}).

Deliverable~3 turned these results into a tool a buyer can use. Every quote is
compared with an expected five-day price range taken from the 10th, 50th and
90th percentiles of past price changes, which works for all four metals; for
gold, the model's predicted direction decides quotes inside the range
(Section~\ref{sec:recommendation}). Because the model signal is weak, the
range does most of the work and the model only tips the close cases. A local
Streamlit application (Section~\ref{sec:application}) runs this inference in
real time on built-in, uploaded or freshly downloaded prices, and explains
each recommendation with the market price, the expected range and the model
probabilities.

\subsection{Future Work}

\begin{itemize}
  \item \textbf{Shift-resistant features}: replace raw price levels with
        relative measures, such as the price relative to its 30-day average,
        and retrain the models on a rolling window so they adapt to new price
        levels.
  \item \textbf{Calibrated decisions}: calibrate the predicted probabilities
        and choose the decision threshold on the validation split instead of
        using 0.5.
  \item \textbf{Forecasts with uncertainty}: Deliverable~1 planned regression
        forecasts with prediction intervals (ARIMA, Prophet and quantile
        models). Deliverable~3 used an empirical range from past price changes
        instead; a fitted forecasting model could give ranges that widen and
        narrow with market volatility.
  \item \textbf{AI models for the other metals}: train and evaluate direction
        models for silver, copper and platinum, so they no longer rely on the
        price range alone.
  \item \textbf{More realistic quotes}: add supplier premiums, local
        currencies and units, and a mode for businesses that sell metal.
\end{itemize}
